\section{Extension and completion in algebraic dimension zero}
\label{sec:completion}

We now combine the preceding constructions. The point of using a
\emph{reduced} boundary is that its whole floor has a section by
Proposition~\ref{floor:section}, while the zero-Lelong metric makes
the obstruction to extending that section vanish if the ambient
adjoint has no sections. We first isolate the elementary finiteness
argument needed for this extension.

\subsection{Finite divisorial support and eventual vanishing}

\begin{lemma}\label{end:finite-primes}
Let \(M\) be a smooth compact K\"ahler manifold with
\(a(M)=q(M)=0\). Then \(M\) has finitely many prime divisors, and
their cohomology classes are linearly independent over \(\R\).
\end{lemma}

\begin{proof}
Suppose a finite collection of distinct prime divisors satisfies a
nonzero rational relation in \(H^2(M,\Q)\). Clearing denominators
and then torsion gives a nonzero integral divisor \(E\) with
\(c_1(\OO_M(E))=0\) in integral cohomology. The exponential
sequence and \(H^1(M,\OO_M)=0\) imply
\(\OO_M(E)\simeq\OO_M\). Its canonical meromorphic section is
then a meromorphic function with divisor \(E\). Since \(a(M)=0\),
every meromorphic function is constant, so \(E=0\), a
contradiction. The classes are rational, so a real dependence would
give a rational dependence by linear algebra. They are therefore
linearly independent over \(\R\); their number is bounded by
\(\dim H^2(M,\R)\).
\end{proof}

\begin{lemma}[Eventual vanishing]\label{end:eventual}
Let \(M\) be as in Lemma~\ref{end:finite-primes}, and let \(L\)
be a rational holomorphic line bundle with \(\kappa(M,L)=-\infty\).
For every fixed holomorphic vector bundle \(\mathcal E\),
\begin{equation}\label{end:fixed-bundle-vanishing}
 H^0\bigl(M,\mathcal E\otimes\OO_M(mL)\bigr)=0
\end{equation}
for all sufficiently large \(m\) in any fixed sufficiently divisible
sequence of integral multiples.
\end{lemma}

\begin{proof}
Assume the contrary. Interpret a section as a morphism
\(\OO_M(-mL)\to\mathcal E\); this requires no meromorphic frame
of \(L\). Among all such morphisms choose a tuple
\(s_1,\ldots,s_k\), with respective indices \(m_1,\ldots,m_k\),
whose generic rank is maximal. Let \(\mathcal H\subset\mathcal E\)
be the saturation of the image of their direct sum. Every further
such morphism has image in \(\mathcal H\): otherwise adding it
would increase generic rank. A nonzero further section can replace
at least one of the \(s_j\) while preserving generic rank.

There are infinitely many increasing indices \(m\) with a nonzero
section, so one fixed replacement index \(j\) works for infinitely
many of them. Taking the nonzero determinant gives effective integral
divisors in the actual line bundles
\begin{equation}\label{end:determinants}
 \det\mathcal H\otimes\OO_M((c+m)L),
 \qquad c=\sum_{i\ne j}m_i.
\end{equation}
The determinant is interpreted reflexively; on the smooth space its
double dual is a line bundle, and the wedge morphism extends to it.

By Lemma~\ref{end:finite-primes}, all the effective divisors in
\eqref{end:determinants} have the same finite set of possible prime
components. Their coefficient vectors lie in \(\Z_{\ge0}^{N}\).
Such an infinite sequence has an infinite componentwise
nondecreasing subsequence: successively pass to a constant or
increasing subsequence in each coordinate. Choose two members with
indices \(m'<m''\). Subtracting their divisors gives an effective
divisor in the actual line \((m''-m')L\), contrary to
\(\kappa(M,L)=-\infty\). This determinant argument is related to
the nonvanishing method in \cite{HLLNonvanishing}.
\end{proof}

\subsection{Extension from a nonempty reduced boundary}

Let \(T_0\) be a smooth compact K\"ahler non-uniruled fourfold
with \(a(T_0)=q(T_0)=0\), and let \(G_0\) be a reduced SNC
divisor. Apply Proposition~\ref{bir:reduced-model}. Write its
nonextracting output as \((T,G)\), with
\[
 A=K_T+G,
 \qquad h:(V,D)\longrightarrow(T,G),
 \qquad J=K_V+D=h^*A.
\]
Here \(T\) is globally strongly \(\Q\)-factorial and klt,
\(K_T\) remains pseudo-effective, \(A\) is analytically nef,
and \((V,D)\) is the projective crepant dlt auxiliary model with
reduced boundary. The pullback of \(J\) to a smooth resolution
has a semipositive singular metric with zero Lelong numbers.

\begin{proposition}\label{end:boundary-nonvanishing}
In this situation, \(G\ne0\) implies \(\kappa(T,A)\ge0\).
\end{proposition}

\begin{proof}
If \(G\ne0\), then \(D\ne0\). By
Proposition~\ref{floor:section}, a positive Cartier multiple of
\(J|_D\) has a nonzero section on the whole reduced divisor
\(D\). Increase this multiple to clear the ambient index. We show
that, under the contrary assumption \(\kappa(T,A)=-\infty\),
high powers of this section extend to \(V\).

Take a simultaneous projective compact K\"ahler log resolution
\(p:M\to V\) also resolving the map to \((T,G)\), and put
\begin{equation}\label{end:crepant}
 L=p^*J,\qquad K_M+F\sim_{\Q}L.
\end{equation}
The support of \(F\) is SNC and every coefficient is at most one;
negative coefficients are allowed. All comparisons here are actual
rational line-bundle comparisons. Normality and projection give
\(\kappa(M,L)=\kappa(V,J)=\kappa(T,A)\). The space \(M\)
still has \(a=q=0\).

For every sufficiently divisible \(m\),
\begin{equation}\label{end:ideal-push}
 p_*\OO_M(mL-\lfloor F\rfloor)
   =\mathcal I_D\otimes\OO_V(mJ).
\end{equation}
Indeed, the nonexceptional coefficient-one components of \(F\)
are the strict transforms of the primes of \(D\). Every other
coefficient-one component has center in \(D\), since \((V,D)\)
is dlt and is klt away from \(D\). A holomorphic function vanishing
on the reduced \(D\) pulls back with order at least one along all
of these components. Conversely, any poles allowed by negative
coefficients of \(\lfloor F\rfloor\) are exceptional. Normality
and Hartogs extension remove them downstairs, while vanishing is
required along each prime of \(D\). This proves
\eqref{end:ideal-push}, after the actual pullback line is removed
by projection.

The low-degree Leray sequence therefore gives an injection
\begin{equation}\label{end:leray}
 H^1\bigl(V,\mathcal I_D\otimes\OO_V(mJ)\bigr)
 \hookrightarrow
 H^1\bigl(M,\OO_M(mL-\lfloor F\rfloor)\bigr).
\end{equation}
No vanishing of a higher direct image is needed for this injection.
Write the line on the right as \(K_M+\mathcal L_m\), where
\begin{equation}\label{end:hl-line}
 \mathcal L_m=
 \OO_M(-K_M-\lfloor F\rfloor+mL)
 \sim_{\Q}(m-1)L+\{F\}.
\end{equation}
For \(m\ge1\), give \(\mathcal L_m\) the metric obtained from
the zero-Lelong semipositive metric on \(L\) and the divisor metric
of \(\{F\}\), taking roots of identified powers if necessary.
Its curvature is semipositive and its multiplier ideal is trivial.
To see the latter directly, the zero-Lelong weight has every finite
local exponential integrability exponent by Skoda's theorem. The
SNC coefficients of \(\{F\}\) are strictly below one, and hence
their divisor weight has an integrability margin. H\"older's
inequality combines that margin with a sufficiently high finite
exponent for the zero-Lelong weight.

Hard Lefschetz with multiplier ideals
\cite[Theorem~0.1]{DPSHardLefschetz} gives a surjection
\[
 H^0(M,\Omega_M^3\otimes\mathcal L_m)
   \longrightarrow H^1(M,K_M+\mathcal L_m).
\]
Its source vanishes for all large divisible \(m\) by
Lemma~\ref{end:eventual}, applied to the fixed bundle
\(\Omega_M^3\otimes\OO_M(-K_M-\lfloor F\rfloor)\).
It follows from \eqref{end:leray} that
\[
 H^1\bigl(V,\mathcal I_D\otimes\OO_V(mJ)\bigr)=0.
\]
The exact sequence for the reduced divisor \(D\) now makes
restriction of sections of \(\OO_V(mJ)\) onto \(D\)
surjective. High divisible powers of the nonzero floor section
remain nonzero because \(D\) is reduced, and therefore extend.
Their extensions push to \(T\), since \(J=h^*A\) and
\(h_*\OO_V=\OO_T\). This contradicts
\(\kappa(T,A)=-\infty\).
\end{proof}

\subsection{Canonical nonvanishing}

\begin{theorem}\label{end:canonical}
Assume Assumptions~\ref{hyp:campana}, \ref{hyp:mmp}, and
\ref{hyp:effective-abundance}. If \(W\) is a smooth compact
K\"ahler non-uniruled fourfold with \(a(W)=0\), then
\(\kappa(W,K_W)\ge0\).
\end{theorem}

\begin{proof}
Suppose \(\kappa(W,K_W)=-\infty\). The Albanese argument
of Lemma~\ref{red:irregular} gives \(q(W)=0\). Both
\(a=q=0\) persist on every smooth compact K\"ahler model used
below. Their canonical classes are pseudo-effective by
non-uniruledness \cite{OuUniruledness}. We distinguish the existence
of a nonzero \emph{meromorphic} section of a positive canonical
power. That property is birationally invariant on smooth spaces:
Jacobian comparison gives it on a common resolution, and
meromorphic tensors extend over codimension two.

\paragraph{No signed canonical frame.}
Suppose first that no positive power of \(K_W\) has a nonzero
meromorphic section. Resolve the finitely many prime divisors on
\(W\), and on the resulting smooth model \(T_0\) let \(G_0\)
be the reduced union of their strict transforms and all exceptional
divisors, with SNC support. This union contains every prime of
\(T_0\): each such prime is either exceptional or maps to a prime
on \(W\). The union may be empty.

Apply Proposition~\ref{bir:reduced-model}. A section of a positive
multiple of \(A=K_T+G\) would, after pullback to a resolution and
division by the meromorphic divisor factors, give a signed
meromorphic pluricanonical tensor. Here \(G\) is rational Cartier
by the global strong condition, and the canonical comparison has a
signed rational exceptional divisor. This is impossible. By
Proposition~\ref{end:boundary-nonvanishing}, it follows that
\(G=0\).

The space \(T\) has no prime divisors: its map from \(T_0\)
extracts none, and the full prime support was included in \(G_0\).
It is klt and has nef \(K_T=A\). Its smooth resolution has
algebraic dimension zero, pseudo-effective canonical class, and no
meromorphic pluricanonical tensor. These are precisely the
hypotheses of Proposition~\ref{fol:empty-case}, which excludes this case.

\paragraph{A signed canonical frame.}
Otherwise resolve the signed divisor of a meromorphic section of
\(mK_W\). On the resulting smooth model write the actual identity
\begin{equation}\label{end:signed}
 K_{T_0}\sim_{\Q}P-N,
\end{equation}
where \(P,N\ge0\) have SNC total support and no common prime.
The pullback tensor includes the resolution discrepancy, so
\eqref{end:signed} is an actual canonical identity. Set
\(G_0=(\Supp P)_{\mathrm{red}}\), and again apply
Proposition~\ref{bir:reduced-model}. Nonextraction and reflexive
extension preserve
\[
 K_T\sim_{\Q}P_T-N_T,
 \qquad G=(\Supp P_T)_{\mathrm{red}}.
\]
If \(G=0\), then \(P_T=0\), and pseudo-effectivity of \(K_T\)
forces \(N_T=0\). Indeed, after pullback to a resolution a
nonzero effective rational Cartier divisor has strictly positive
K\"ahler mass, whereas its negative cannot represent a
pseudo-effective class.

If \(G\ne0\), Proposition~\ref{end:boundary-nonvanishing}
gives a section of a multiple of \(A\). The signed divisor
\(P_T+G-N_T\) representing \(A\) is unique: two meromorphic
sections of the same power differ by a meromorphic function, and
\(a(T)=0\). The section's divisor is effective. Since \(P_T+G\)
and \(N_T\) have disjoint support, we again obtain \(N_T=0\).
Thus in both cases \(K_T\) has nonvanishing.

Run Assumption~\ref{hyp:mmp} on the exact ordinary klt pair
\((T,0)\), and denote its nef endpoint by \(T'\). Its input
is globally strongly \(\Q\)-factorial and its canonical class
is pseudo-effective, as required. Canonical sections push forward
through the nonextracting program. Applying
Assumption~\ref{hyp:effective-abundance} therefore makes
\(K_{T'}\) semiample. Algebraic dimension zero forces the resulting
map to have a point image, so \(K_{T'}\) is actually torsion.

It remains to return to a \emph{smooth} canonical bundle. On a
smooth compact K\"ahler resolution \(\mu:M'\to T'\),
\begin{equation}\label{end:exceptional}
 K_{M'}\sim_{\Q}\mu^*K_{T'}+E\sim_{\Q}E
\end{equation}
for a possibly signed exceptional rational divisor \(E\). Let
\(\Theta\) be a positive current in the pseudo-effective class
\(c_1(K_{M'})\). If \(\omega_{T'}\) is a K\"ahler form on
\(T'\), exceptionality gives
\[
 \int_{M'}\Theta\wedge\mu^*\omega_{T'}^3=0.
\]
The measure is nonnegative and \(\mu^*\omega_{T'}\) is strictly
positive on the isomorphism locus. Hence \(\Theta\) is supported
in the analytic exceptional/non-isomorphism locus. The support
theorem for positive closed \((1,1)\)-currents expresses it as an
effective real sum of divisorial currents
\cite{DemaillyComplexAnalyticGeometry}. By
Lemma~\ref{end:finite-primes}, the prime classes on \(M'\) are
linearly independent. Comparison with \eqref{end:exceptional}
therefore forces every coefficient of \(E\) to be nonnegative.
The actual identity \(K_{M'}\sim_{\Q}E\ge0\) contradicts
smooth birational invariance of Kodaira dimension and completes
the proof.
\end{proof}

\subsection{The good minimal model}

\begin{proof}[Proof of Theorems~\ref{thm:nonvanishing} and
\ref{thm:main}]
The reduction in Proposition~\ref{red:reduction} treats the
projective, uniruled and irregular cases and leaves smooth
non-uniruled non-Moishezon fourfolds with \(q=0\). Positive
algebraic dimension is handled by Proposition~\ref{pos:nonvanishing};
algebraic dimension zero is handled by
Theorem~\ref{end:canonical}. The sections push to the actual
adjoint of the nef pair exactly as in the reduction. This proves
Theorem~\ref{thm:nonvanishing}.

Finally start Assumption~\ref{hyp:mmp} from the original pair
\((X,B)\). Its endpoint \((Y,B_Y)\) is in the required global
strong category, with nef actual adjoint, no extraction, and all the
stated discrepancy inequalities. Theorem~\ref{thm:nonvanishing}
gives a first section there. Assumption~\ref{hyp:effective-abundance}
then makes a positive Cartier multiple of that \emph{same endpoint}
adjoint globally generated. These are all the assertions of
Theorem~\ref{thm:main}.
\end{proof}
